\appendix
\renewcommand{\appendixtocname}{Appendix}
\renewcommand{\appendixpagename}{\appendixtocname}
\addappheadtotoc
\appendixpage

% ═════════════════════════════════════════════════════════════════════════════
\chapter{API Reference}\label{app:api}
% ═════════════════════════════════════════════════════════════════════════════

Table~\ref{tab:api_controllers} lists all REST API controllers, their base routes,
and the primary endpoints they expose.
All routes are prefixed with \texttt{/api}.
Every endpoint except the authentication endpoints and the connectivity check requires
a valid \texttt{X-User-Email} header, resolved by the local-identity middleware
described in Section~\ref{sec:meth_backend}.

\begin{table}[H]
\centering
\caption{REST API controllers and primary endpoints.}
\label{tab:api_controllers}
\footnotesize
\begin{tabular}{>{\raggedright\arraybackslash}p{4.0cm}>{\raggedright\arraybackslash}p{4.0cm}>{\raggedright\arraybackslash}p{6.0cm}}
\toprule
\textbf{Controller} & \textbf{Base route} & \textbf{Key endpoints} \\
\midrule
AuthController
  & \texttt{/api/auth}
  & POST \texttt{/register}, POST \texttt{/login},
    POST \texttt{/forgot-password}, POST \texttt{/reset-password} \\
\addlinespace
TasksController
  & \texttt{/api/tasks}
  & GET, POST \texttt{/},
    GET, PUT, DELETE \texttt{/\{id\}},
    PATCH \texttt{/\{id\}/status},
    PATCH \texttt{/\{id\}/star},
    GET \texttt{/search} \\
\addlinespace
ProjectsController
  & \texttt{/api/projects}
  & GET, POST \texttt{/},
    GET, PUT, DELETE \texttt{/\{id\}},
    POST \texttt{/\{id\}/invite},
    GET \texttt{/\{id\}/members} \\
\addlinespace
TeamsController
  & \texttt{/api/teams}
  & GET, POST \texttt{/},
    GET, PUT, DELETE \texttt{/\{id\}},
    POST \texttt{/\{id\}/invite},
    POST \texttt{/invitations/\{id\}/accept} \\
\addlinespace
MessagesController
  & \texttt{/api/messages}
  & GET \texttt{/conversations},
    GET \texttt{/\{userId\}},
    POST \texttt{/},
    DELETE \texttt{/conversation/\{userId\}} \\
\addlinespace
Group\-Chats\-Controller
  & \texttt{/api/groupchats}
  & GET, POST \texttt{/},
    GET, PUT, DELETE \texttt{/\{id\}},
    POST \texttt{/\{id\}/members},
    POST \texttt{/\{id\}/messages} \\
\addlinespace
Notifi\-cations\-Controller
  & \texttt{/api/}\cbrk\texttt{notifications}
  & GET \texttt{/},
    GET \texttt{/unread-count},
    PATCH \texttt{/\{id\}/read},
    PATCH \texttt{/read-all},
    DELETE \texttt{/\{id\}} \\
\addlinespace
Calendar\-Events\-Controller
  & \texttt{/api/}\cbrk\texttt{calendarevents}
  & GET, POST \texttt{/},
    GET, PUT, DELETE \texttt{/\{id\}} \\
\addlinespace
Chatbot\-Controller
  & \texttt{/api/chatbot}
  & GET \texttt{/conversations},
    POST \texttt{/conversations},
    GET \texttt{/conversations/}\cbrk\texttt{\{id\}/messages},
    POST \texttt{/conversations/}\cbrk\texttt{\{id\}/messages},
    POST \texttt{/scan} \\
\addlinespace
Dash\-board\-Controller
  & \texttt{/api/dashboard}
  & GET \texttt{/stats},
    GET \texttt{/due-this-week},
    GET \texttt{/recent-activity} \\
\addlinespace
Task\-Comments\-Controller
  & \texttt{/api/}\cbrk\texttt{tasks/}\cbrk\texttt{\{taskId\}/comments}
  & GET, POST \texttt{/},
    DELETE \texttt{/\{id\}} \\
\addlinespace
SettingsController
  & \texttt{/api/settings}
  & GET, PUT \texttt{/profile},
    PUT \texttt{/password} \\
\addlinespace
Connec\-tivity\-Controller
  & \texttt{/api/}\cbrk\texttt{connectivity}
  & GET \texttt{/ping},
    POST \texttt{/bulk-sync} \\
\addlinespace
DevController
  & \texttt{/api/dev}
  & GET \texttt{/db-info}
    (development environment only) \\
\bottomrule
\end{tabular}
\end{table}

% ═════════════════════════════════════════════════════════════════════════════
\chapter{Database Schema}\label{app:schema}
% ═════════════════════════════════════════════════════════════════════════════

\section*{SQLite Entities}

Table~\ref{tab:sqlite_entities} lists all eighteen SQLite entities managed by Entity
Framework Core, together with their primary key type, key fields, and any sync
participation.

\begin{table}[H]
\centering
\caption{SQLite entity summary.}
\label{tab:sqlite_entities}
\footnotesize
\begin{tabular}{>{\raggedright\arraybackslash}p{3.5cm}l>{\raggedright\arraybackslash}p{8.0cm}}
\toprule
\textbf{Entity} & \textbf{PK} & \textbf{Key fields} \\
\midrule
\texttt{AppUser}
  & \texttt{int Id}
  & FirstName, LastName, Email, PasswordHash, AvatarUrl, Company,
    Country, Phone, Timezone, ResetToken, ResetTokenExpiry \\
\addlinespace
\texttt{TaskItem}
  & \texttt{int Id}
  & Title, Description, Status, Priority, DueDate, IsStarred,
    ProjectId, AssigneeId, SyncId, UpdatedAt, IsSynced \\
\addlinespace
\texttt{Project}
  & \texttt{int Id}
  & Name, Description, Color, OwnerId, IsStarred,
    SyncId, UpdatedAt, IsSynced \\
\addlinespace
\texttt{ProjectMember}
  & \texttt{int Id}
  & ProjectId, UserId, Role, JoinedAt \\
\addlinespace
\texttt{Team}
  & \texttt{int Id}
  & Name, Description, OwnerId, CreatedAt \\
\addlinespace
\texttt{TeamMember}
  & \texttt{int Id}
  & TeamId, UserId, Role, JoinedAt \\
\addlinespace
\texttt{Message}
  & \texttt{int Id}
  & SenderId, ReceiverId, Body, IsRead, SentAt,
    AttachmentUrl, IsDeletedBySender, IsDeletedByReceiver, IsEdited \\
\addlinespace
\texttt{GroupChat}
  & \texttt{int Id}
  & Name, CreatedByUserId, CreatedAt;
    child: \texttt{GroupChatMember} (GroupChatId, UserId, JoinedAt);
    child: \texttt{GroupChatMessage} (GroupChatId, SenderId, Body) \\
\addlinespace
\texttt{Notification}
  & \texttt{int Id}
  & UserId, Title, Message, NotificationType, IsRead,
    ActionUrl, RelatedTaskId, SyncId, UpdatedAt \\
\addlinespace
\texttt{CalendarEvent}
  & \texttt{int Id}
  & OwnerId, Title, Description, StartAt, EndAt,
    Color, MeetingLink \\
\addlinespace
\texttt{Chatbot}\cbrk\texttt{Conversation}
  & \texttt{int Id}
  & UserId, Title, Mode, CreatedAt, UpdatedAt \\
\addlinespace
\texttt{ChatbotMessage}
  & \texttt{int Id}
  & ConversationId, Role, Content, CreatedAt \\
\addlinespace
\texttt{Reminder}
  & \texttt{int Id}
  & TaskId, UserId, FireAt, NotifyEmail, NotifyInApp,
    HasFired, FiredAt, SyncId \\
\addlinespace
\texttt{SyncOutboxEntry}
  & \texttt{int Id}
  & OperationId, OperationName, PayloadJson, Status
    (Pending/Processing/Synced/Failed), Attempts, LastAttemptAt, ErrorMessage \\
\addlinespace
\texttt{TaskComment}
  & \texttt{int Id}
  & TaskId, AuthorId, Body, CreatedAt \\
\addlinespace
\texttt{LocalInvitation}
  & \texttt{int Id}
  & InviteeId, InviterId, TeamId or ProjectId, Status, CreatedAt \\
\addlinespace
\texttt{LocalTeamMember}
  & \texttt{int Id}
  & TeamId, UserId, Role, JoinedAt (local mirror of MongoDB record) \\
\addlinespace
\texttt{ISyncableEntity}
  & ---
  & Interface only: SyncId, UpdatedAt, IsSynced, ToMongoDocument() \\
\bottomrule
\end{tabular}
\end{table}

\section*{MongoDB Collections}

Table~\ref{tab:mongo_appendix} lists the five MongoDB collections used as the relay
layer.

\begin{table}[H]
\centering
\caption{MongoDB collection descriptions.}
\label{tab:mongo_appendix}
\begin{tabular}{p{3.5cm}p{9.5cm}}
\toprule
\textbf{Collection} & \textbf{Description and key fields} \\
\midrule
\texttt{tasks}
  & Mirror of \texttt{TaskItem}. Documents keyed by \texttt{syncId} (GUID).
    Fields: syncId, title, description, status, priority, dueDate, ownerId,
    projectId, isStarred, updatedAt. \\
\addlinespace
\texttt{notifications}
  & Mirror of \texttt{Notification}. Fields: syncId, userId, title, message,
    notificationType, isRead, actionUrl, relatedTaskId, updatedAt. \\
\addlinespace
\texttt{team\_}\cbrk\texttt{invitations}
  & Cross-device invitation relay. Fields: inviterEmail, inviteeEmail,
    teamId, projectId, teamName, status (Pending/Accepted/Declined), createdAt. \\
\addlinespace
\texttt{team\_members}
  & Mirror of accepted team memberships. Fields: teamId, userId,
    userEmail, role, joinedAt. \\
\addlinespace
\texttt{user\_presence}
  & Optional online-status documents. Fields: userId, lastSeenAt, isOnline. \\
\bottomrule
\end{tabular}
\end{table}

% ═════════════════════════════════════════════════════════════════════════════
\chapter{Configuration Reference}\label{app:config}
% ═════════════════════════════════════════════════════════════════════════════

\section*{appsettings.json Structure}

The application is configured via \texttt{appsettings.json}
(overridden by \texttt{appsettings.Development.json} in development).
Table~\ref{tab:appsettings} describes each configuration section.

\enlargethispage{\baselineskip}
\begin{table}[H]
\centering
\small
\caption{\texttt{appsettings.json} section reference.}
\label{tab:appsettings}
\begin{tabular}{>{\raggedright\arraybackslash}p{5.5cm}>{\raggedright\arraybackslash}p{7.5cm}}
\toprule
\textbf{Section / key} & \textbf{Description} \\
\midrule
\texttt{DatabaseProvider}
  & \texttt{"SQLite"} (default) or \texttt{"SqlServer"}.
    Selects the EF Core provider at startup. \\
\addlinespace
\texttt{ConnectionStrings.}\cbrk\texttt{DefaultConnection}
  & SQLite connection string, e.g.\ \texttt{Data Source=taskflow.db}.
    Relative paths are resolved against the executable directory. \\
\addlinespace
\texttt{MongoDB.}\cbrk\texttt{ConnectionString}
  & Full MongoDB connection string including credentials.
    Leave empty to run in SQLite-only mode without sync. \\
\addlinespace
\texttt{Jwt.Key}
  & HMAC-SHA256 signing key (minimum 32 characters) \emph{(vestigial ---
    the local-identity middleware does not issue or validate JWT tokens)}. \\
\addlinespace
\texttt{Jwt.Issuer} / \texttt{.Audience}
  & Token issuer and audience strings \emph{(vestigial)}. \\
\addlinespace
\texttt{Jwt.}\cbrk\texttt{ExpiresInMinutes}
  & Token lifetime in minutes \emph{(vestigial)}. \\
\addlinespace
\texttt{Mistral.Model}
  & Model ID for the General chatbot mode.
    Default: \texttt{mistral-small-latest}. \\
\addlinespace
\texttt{Mistral.}\cbrk\texttt{CoderModel}
  & Model ID for the Coder mode.
    Default: \texttt{codestral-latest}. \\
\addlinespace
\texttt{Mistral.}\cbrk\texttt{OcrModel}
  & Model ID for the Scanner (OCR) mode.
    Default: \texttt{mistral-ocr-latest}. \\
\addlinespace
\texttt{Mistral.}\cbrk\texttt{HistoryWindow}
  & Number of prior messages included in the Mistral context.
    Default: 20. \\
\addlinespace
\texttt{EmailSettings.*}
  & SMTP configuration for MailKit.
    SmtpHost, SmtpPort (587), UseSsl (false for STARTTLS),
    SenderEmail, SenderName, Username, Password. \\
\addlinespace
\texttt{Kestrel.}\cbrk\texttt{Endpoints.}\cbrk\texttt{Http.Url}
  & \texttt{http://127.0.0.1:0} — binds to a random available port.
    The chosen port is printed to stdout as
    \texttt{TASKFLOW\_}\cbrk\texttt{BACKEND\_READY:\{address\}}. \\
\bottomrule
\end{tabular}
\end{table}

\section*{Environment Variables}

Table~\ref{tab:env_vars} lists the environment variables recognised by the application
at startup.

\begin{table}[H]
\centering
\caption{Supported environment variables.}
\label{tab:env_vars}
\begin{tabular}{>{\raggedright\arraybackslash}p{5.2cm}>{\raggedright\arraybackslash}p{7.8cm}}
\toprule
\textbf{Variable} & \textbf{Effect} \\
\midrule
\texttt{ASPNETCORE\_}\cbrk\texttt{ENVIRONMENT}
  & Standard ASP.NET Core environment selector.
    Set to \texttt{Development} to enable Swagger UI and verbose logging. \\
\addlinespace
\texttt{ASPNETCORE\_URLS}
  & Overrides Kestrel endpoint binding when set.
    Normally not set; the \texttt{Kestrel} config section is used instead. \\
\bottomrule
\end{tabular}
\end{table}

\section*{tauri.conf.json Summary}

Table~\ref{tab:tauri_config} summarises the key \texttt{tauri.conf.json}
settings for the Tauri-based installer.

\begin{table}[H]
\centering
\caption{Key \texttt{tauri.conf.json} configuration (Tauri build).}
\label{tab:tauri_config}
\begin{tabular}{>{\raggedright\arraybackslash}p{5.5cm}>{\raggedright\arraybackslash}p{7.5cm}}
\toprule
\textbf{Key} & \textbf{Value / description} \\
\midrule
\texttt{productName}  & \texttt{Task Flow} \\
\texttt{identifier}   & \texttt{com.taskflow.desktop} \\
\texttt{build.}\cbrk\texttt{frontendDist} & \texttt{../wwwroot} — directory containing \texttt{index.html}
                              and Webpack output bundles. \\
\texttt{app.}\cbrk\texttt{windows[0]} & Main window: $1280\times800$, centred, initially hidden,
                          show on backend-ready signal. \\
\texttt{app.}\cbrk\texttt{windows[1]} & Splashscreen: $400\times300$, undecorated, initially hidden,
                          shown during backend startup. \\
\texttt{app.}\cbrk\texttt{security.csp} & Content Security Policy allowing localhost API and
                            SignalR WebSocket connections. \\
\texttt{bundle.}\cbrk\texttt{targets}
  & \texttt{["nsis", "msi"]} — Windows installer formats. \\
\texttt{bundle.}\cbrk\texttt{externalBin}
  & \texttt{["binaries/taskflow-backend"]} — sidecar path. \\
\texttt{bundle.windows.}\cbrk\texttt{webview}\cbrk\texttt{InstallMode}
  & \texttt{downloadBootstrapper} — download WebView2 runtime if absent. \\
\texttt{bundle.windows.}\cbrk\texttt{nsis.}\cbrk\texttt{installMode}
  & \texttt{currentUser} — no administrator privileges required. \\
\texttt{plugins.}\cbrk\texttt{shell.open}
  & \texttt{true} — permits opening URLs in the default browser. \\
\bottomrule
\end{tabular}
\end{table}
